% Unregistered auxiliary research. Complete proof supplied for independent audit.
\paragraph{Supported higher families (auxiliary proof under review).}
Let $k$ be algebraically closed of characteristic $p>0$ and let
$E=\bigcup_iE_i$ be a projective SNC surface with smooth connected
components. No negative exceptional divisor is needed in this lemma.
Let $u\in k^*$ and choose a finite tuple
\[
             z_i\in\operatorname{CH}_0(E_i)^0_{\mathbf Q}.
\]
Regard these as their canonical top classes in $K_0(E_i)_{\mathbf Q}$.
Then there exist a smooth connected projective curve $P$, points
$o,t\in P(k)$, and an actual class
$\eta\in K_1(E\times P)_{\mathbf Q}$ satisfying
\begin{equation}
 \bigl((\eta_t-\eta_o)|_{E_i}\bigr)_i=(u z_i)_i
                  \quad\hbox{in }\bigoplus_iK_1(E_i)_{\mathbf Q}.
                                                        \tag{HK.1}
\end{equation}
In fact the two component endpoint tuples can be chosen to be
$(0)_i$ and $(u z_i)_i$ separately. The theorem does not identify
the global endpoint difference with $u$ times a selected global
$K_0(E)$ lift. It makes no assertion about a proper relation in
$K_0(E)$ or about the quotient (UK.7).

\emph{1. Represent the component cycles on transverse proper curves.}
Clear one common denominator. Move each integral degree-zero
representative away from all double curves, using the explicit
smooth-curve moving argument of C.3br Step~4. Choose a smooth
connected projective curve $A_i\subset E_i$ through its finitely
many support points, transverse to every double curve and avoiding
all triple points. A sufficiently high very ample system separates
these finite values and first jets. The same incidence dimension
counts and connectedness argument as in C.3br give the required
curve in every characteristic. Put
\[
 L_i=\mathcal O_{A_i}\left(\sum_x n_{i,x}[x]\right)
                   \in\operatorname{Pic}^0(A_i)(k).
\]
For $a_i:A_i\hookrightarrow E_i$ the exact sequences of point
Cartier divisors on the smooth curve give the integral identity
\begin{equation}
              (a_i)_*([L_i]-1)=\sum_x n_{i,x}[\mathcal O_x]
                      \quad\hbox{in }K_0(E_i).          \tag{HK.2}
\end{equation}
Thus no rational cohomological representative has replaced the
specified zero-cycle class.

\emph{2. Use the complete marked edge-family lemma.}
Apply the integral lemma (VE.2) simultaneously to the curves $A_i$
with line bundles $L_i$ and the fixed unit $u$. Mark every point
of $A_i$ on every double curve of E. It yields one actual smooth
connected proper parameter $P$ and classes
\[
 \xi_i\in K_1(A_i\times P),\qquad
  (\xi_i)_o=0,\qquad(\xi_i)_t=u([L_i]-1),
 \qquad \xi_i|_{\{q\}\times P}=0
\]
for every marking q. The final equality holds in the actual
$K_1(P)$, not just at the endpoints or after a norm.

\emph{3. Retain supports to remove all triple coherence.}
Let $W\subset E$ be the reduced closed union of the finitely many
$A_i$. It avoids every triple point. Write $W_i=W\cap E_i$ and
$W_e=W\cap C_e$; each $W_e$ is a finite reduced set.
Use nonconnective perfect K-theory with the indicated support.
The regular closed immersion $A_i\times P\hookrightarrow E_i\times P$
gives a proper perfect pushforward to the supported group
\[
 \beta_i=(a_i\times1)_*\xi_i
       \in K_1(E_i\times P\ \mathrm{on}\ W_i\times P)_{\mathbf Q}.
\]
The support is enlarged from $A_i\times P$ to $W_i\times P$ if
necessary. Perfect proper pushforward and its supported form are
\citet[Sections 3.16.4--3.16.7]{D_ThomasonTrobaugh1990}.

For an incident edge $C_e\subset E_i$, transversality makes the
square with $A_i\subset E_i$ Tor-independent. Its derived base-change
homotopy, including supports, identifies the restriction of $\beta_i$
with the sum of the pushforwards of
$\xi_i|_{\{q\}\times P}$ for $q\in A_i\cap C_e$.
Every one is zero. This is zero in
$K_1(C_e\times P\ \mathrm{on}\ W_e\times P)$ itself.
If $W_e$ includes further points arising from a neighbouring
component, enlarging support preserves zero. The exact primary
base-change statement and full proof are
\citet[Proposition 3.18, pp.~321--322]{D_ThomasonTrobaugh1990}.
Consequently the alternating edge restriction of $(\beta_i)_i$ is zero.

Apply the rational closed-union descent of (ET.6)--(UK.2) both to
$E\times P$ and to its complement of $W\times P$.
Finite homotopy limits commute with the fibre defining support
K-theory. Hence the supported descent diagram is the actual fibre
of those two coherent diagrams, with the canonical geometric
nullhomotopies retained. All its vertex support spectra are zero,
because W avoids the vertices. It therefore gives the fibre sequence
\begin{equation}
 K(E\times P\ \mathrm{on}\ W\times P)_{\mathbf Q}
 \longrightarrow
 \bigoplus_iK(E_i\times P\ \mathrm{on}\ W_i\times P)_{\mathbf Q}
 \longrightarrow
 \bigoplus_eK(C_e\times P\ \mathrm{on}\ W_e\times P)_{\mathbf Q}.
                                                        \tag{HK.3}
\end{equation}
Its long exact sequence lifts $(\beta_i)_i$ to an actual supported
$K_1$ class. Forget its support to obtain $\eta$.
There is no use of a perfect pushforward $K(A_i)\to K(E)$, which
would be unjustified at the double points. Only the regular component
pushforwards and actual supported descent are used.

\emph{4. Evaluate the actual proper family.}
Both endpoint squares with $a_i\times1$ are Tor-independent.
Perfect derived base change and the projection formula give
\[
 \eta_o|_{E_i}=0,\qquad
 \eta_t|_{E_i}=u(a_i)_*([L_i]-1)=u z_i.
\]
The complete projection-formula proof is
\citet[Proposition 3.17, pp.~320--321]{D_ThomasonTrobaugh1990}.
Divide by the cleared common denominator. This proves (HK.1).
No Adams decomposition or semisimplicity of $K_1$ has been used.
The construction lives on the fixed proper product $E\times P$.

\emph{Why this is not yet a solution of the component gate.}
The statement raises the K-degree by multiplication with u. A
valuation of u over an auxiliary function field may recover the
component zero-cycle tuple, but compatibility of that residue with
ordinary proper $K_0$ relations is an additional theorem. A regular
semistable parameter model can have a special dual graph with cycles;
its support Mayer--Vietoris boundary has node coefficients in
$K_{-1}(E)$. Their vanishing after point pushforward is not the
vanishing of the abstract Picard tensor in (PL.5).
C.3bt removes global endpoint ambiguities with zero component image,
but does not by itself supply this missing specialization theorem.
The previously retained negative-K valuation counterexample remains
valid and rules out an unrestricted residue principle.
